 \documentclass{resume} % Use the custom resume.cls style
  \usepackage[dvipsnames]{xcolor}
  \usepackage[hidelinks]{hyperref}

  \DeclareTextCommandDefault{\&}{\textnormal{\symbol{38}}}

  \usepackage[left=0.65in,top=0.5in,right=0.65in,bottom=0.4in]{geometry}
  \newcommand{\tab}[1]{\hspace{.2667\textwidth}\rlap{#1}}
  \newcommand{\itab}[1]{\hspace{0em}\rlap{#1}}

  \name{Qiwei Ma}
  \address{Homepage: \href{https://qiwei-ma.github.io/}{\texttt{https://qiwei-ma.github.io/}}}
  \address{Hong Kong, China}
  \address{Phone: (+852) 70115761 \\ Email: \texttt{mqiwei@eduhk.hk}}

  \definecolor{CarnegieMellonRed}{RGB}{196,18,48}

  \renewenvironment{rSection}[1]{
  \sectionskip
  \textbf{\MakeUppercase{#1}}
  \sectionlineskip
  \hrule
  \begin{list}{}{
  \setlength{\leftmargin}{1.5em}
  }
  \item[]
  }{
  \end{list}
  }

  \begin{document}

  \begin{rSection}{Education}
  {\bf Shenzhen University (SZU)} \hfill {\em June 2026} \\
  M.S. in Artificial Intelligence \\
  GPA: 3.5/4.0

  {\bf Ningxia University (NXU)} \hfill {\em June 2022} \\
  B.E. in Materials \\
  GPA: 3.0/4.0
  \end{rSection}

  \begin{rSection}{Interests}
  \begin{rSubsection}{}{}{}{}
  \item Robotics and Embodied Intelligence
  \item AI in Education
  \item Large Language Models and Multi-Agent Systems
  \end{rSubsection}
  \end{rSection}

  \begin{rSection}{Research Experience}

  \begin{rSubsection}{Multi-Agent Conversational AI for EFL Speaking Practice \cite{paper1}}{Jan
  2025 - Jan 2026}{Supervisor: \href{https://fe.szu.edu.cn/info/1023/1133.htm}{Dr. Zhang}}{SZU}
  \item Developed a multi-agent system (MAS) for EFL speaking practice.
  \item Built seven specialized agents for preprocessing, response generation, and dialogue
  supervision.
  \item Investigated the mechanisms underlying MAS superiority and identified synergistic effects
  among integrated features.
  \item Conducted a 4-week controlled experiment with 32 university EFL learners, showing that MAS
  outperformed single-agent systems (SAS) in oral proficiency gains (p = 0.049) and grammatical
  accuracy (p = 0.016).
  \item Found that MAS led to 26\% more practice sessions, 15\% longer utterances, and a 70\%
  reduction in repeated grammatical errors.
  \end{rSubsection}

  \begin{rSubsection}{Intelligent Annotation and Feedback System for English Writing \cite{paper4}}
  {Apr 2025 - Sep 2025}{Supervisor: \href{https://fe.szu.edu.cn/info/1023/1133.htm}{Dr. Zhang}}{SZU}
  \item Developed MLLM-IAF (MLLM-based Intelligent Annotation and Feedback), a mobile
  application for automated English writing assessment for junior high school students.
  \item Implemented a dual-engine AI workflow combining semantic evaluation and visual grounding to
  provide immediate visual feedback with error localization on handwritten essays.
  \item Conducted an experimental study comparing students using MLLM-IAF with a control group
  receiving traditional instruction across four weekly writing tasks.
  \item Observed significant improvements in writing performance, learning engagement, and writing
  self-efficacy in the experimental group.
  \item Demonstrated strong agreement between AI and teacher scoring, and students reported high
  satisfaction with the system's usefulness and ease of use.
  \end{rSubsection}

  \end{rSection}

\clearpage

  \begin{rSection}{Additional Research}

  \begin{rSubsection}{Reasoning for Table Manipulation \textsuperscript{*}}{Mar 2025 - Jan 2026}
  {Supervisors: \href{https://minyang.me/}{Prof. Yang}, \href{https://szs.siat.ac.cn/siat/2025-02/21/article_2025022105340818693.html}{Dr. Tan}}{SIAT, CAS}
  \item Developed an end-to-end LLM framework for table manipulation through structured reasoning.
  \item Constructed a benchmark covering five core tasks: table splitting and merging, wide-to-long
  conversion, semi-structured field parsing, and row/column generation.
  \item Completed two-stage training with supervised fine-tuning on reasoning traces followed by
  GRPO optimization, achieving state-of-the-art performance among 7B-scale table-specific models.
  \item Analyzed structural integrity challenges and found that column-level accuracy often exceeded
  row-level accuracy because of sensitivity to missing fields.
  \end{rSubsection}

  \begin{rSubsection}{Differential Privacy Image Generation}{Mar 2024 - Apr 2025}{Supervisor:
  \href{https://fe.szu.edu.cn/info/1023/1133.htm}{Dr. Zhang}}{SZU}
  \item Developed a differential privacy framework for image generation to reduce gradient clipping
  bias and improve training stability.
  \item Introduced reconstruction loss and noise injection during generator upsampling to improve
  data utility and image diversity.
  \item Designed a multi-component training pipeline integrating a generator, discriminators, a
  classifier, and an encoder with gradient sanitization mechanisms.
  \item Achieved state-of-the-art performance on MNIST and Fashion-MNIST, surpassing baseline
  methods in Inception Score, Frechet Inception Distance, and downstream classification accuracy.
  \end{rSubsection}

  \end{rSection}

  \begin{rSection}{Publications}
  \itemsep -2pt

  \begingroup\small
  \renewcommand{\section}[2]{}
  \begin{thebibliography}{9}
  \setlength{\itemsep}{2pt}
  \setlength{\parskip}{0pt}
  \bibitem{paper1} Jun Zhang, \textbf{Qiwei Ma}, Yu Zhang, and Xiaoming Cao. (2026).
  Multi-agent vs. single-agent AI for EFL speaking practice: A controlled experiment with hybrid input,
  contextual dialogue, and proficiency-adaptive feedback. \emph{Educational Technology \& Society},
  29(2), 297--322. \href{https://doi.org/10.30191/ETS.202604_29(2).SP05}{doi:10.30191/ETS.202604\_29(2).SP05}.

  \bibitem{paper4} Jun Zhang, \textbf{Qiwei Ma}, Yu Zhang, Liying Zhu, and Xiaoming Cao. (2027).
  From textual lists to in-situ annotations: Empowering junior high school EFL writing with multimodal
  large language models. \emph{Assessing Writing}, 71, 101146.
  \href{https://doi.org/10.1016/j.asw.2026.101146}{doi:10.1016/j.asw.2026.101146}.

  \bibitem{tlt2026} Juan Yang, Minjuan Wang, Xu Du, and \textbf{Qiwei Ma}. (2026).
  Multimodal affective computing in the era of LLM-based educational agents: Key insights,
  challenges, and future prospects. \emph{IEEE Transactions on Learning Technologies}, 19, 988--1004.
  \href{https://doi.org/10.1109/TLT.2026.3738338}{doi:10.1109/TLT.2026.3738338}.

  \bibitem{agentic2026} Minjuan Wang, Yizhi Wang, and \textbf{Qiwei Ma}. (2026).
  Agentic AI unleashed: Power, promise, and ethical frontiers. In K. T. Chui, C. Li, J. Jia,
  S. K. S. Cheung, F. W.-F. Lee, and J. Shang (Eds.),
  \emph{Blended learning: Innovations for future education} (pp. 27--39). Springer.
  \href{https://doi.org/10.1007/978-981-92-3761-6_3}{doi:10.1007/978-981-92-3761-6\_3}.
  \end{thebibliography}
  \endgroup
  \end{rSection}


  \begin{rSection}{Additional Manuscript}
  \textsuperscript{*}\textbf{Qiwei Ma}, Bo Li, Hengyang Mu, Chuancheng Zhao, Hong Li, Kai Yan,
  Qingwen Li, Jun Zhang, Minghuan Tan, and Min Yang.
  Advancing large language models with reasoning for dynamic and hierarchical real-world table manipulation.
  \end{rSection}

  \begin{rSection}{Experience}

  \begin{rSubsection}{Research Assistant}{Apr 2026 - Apr 2027}{The Education University of Hong
  Kong}{Hong Kong, China}
  \item Research focus: robotics, embodied intelligence, and AI for education.
  \end{rSubsection}

  \begin{rSubsection}{Research Intern}{May 2025 - Jan 2026}{\href{http://english.siat.cas.cn/}
  {Shenzhen Institutes of Advanced Technology, Chinese Academy of Sciences}}{China}
  \item Worked on table reasoning and manipulation with large language models.
  \end{rSubsection}

  \end{rSection}

  \end{document}